\section{Design Implications and Next Steps}
\subsection{From established paradigms to new combinations}
The first test of a design representation is whether it can express competitive existing implementations. Attention provides a useful progression: materialized scores, tiled fusion, online normalization, different work partitions, and hardware-supported data movement. The representation should expose the mechanisms that distinguish these families, allowing them to combine with hardware choices. Encoding the name of an existing kernel as one opaque option offers a useful baseline, but provides little opportunity to recombine its internal decisions.

A richer representation creates the possibility of improving on established designs; it does not establish that improvement by itself. A frontier comparison needs a current strong implementation for the same workload, precision, and target constraints. Software comparisons hold hardware fixed. System comparisons allow hardware changes under a common resource and implementation-cost model. Reporting both separates algorithmic gains from additional hardware capability. The Ventus studies establish useful combinations within their declared domains; a claim of performance beyond state-of-the-art implementations requires that additional comparison.

\subsection{How much of the boundary should move?}
A software--hardware boundary should be reopened when its fixed assumptions prevent a useful design or impose repeated overhead. Reopening every interface simultaneously makes both implementation and optimization harder. The proposed approach introduces explicit alternatives for a selected mechanism, retains their programmable degrees of freedom, and evaluates the consequences over the required workload set.

The current cases vary hardware resources and software organization within an existing interface. The boundary formulation in Section~\ref{sec:boundary-example} extends that domain. An implementation study of boundary movement should generate the competing controllers or instruction mechanisms, compile matching programs, and include area and timing costs under a common technology and clock methodology. This would test whether specialization's recurring performance benefit justifies its fixed cost and reduced flexibility.

\subsection{Cost fidelity and optimization scale}
Fine-grained timing, bank conflicts, and dynamic arbitration can overwhelm a global MILP. A compact master model with conditional resource signatures and selective detailed completion is one way to manage this scale. It introduces a modeling choice: which interactions can be summarized without changing important rankings? The model-fidelity case shows why this question needs execution evidence.

The quality of the implementation library also limits the domain. Compiler register allocation, unsupported layouts, and backend synchronization rules can change feasibility after an initial model is built. These effects should produce explicit family restrictions or refined estimates. Physical cost introduces additional interactions, including memory macros, wiring, frequency, and control logic. A logical-area estimate and a cycle count do not determine an implemented chip's energy or latency.

\subsection{Evaluating the agent and reusable knowledge}
An agent can improve the design space, the formulation, or the allocation of search effort. These effects require different controls. With a fixed final space, compare equivalent formulations and structured search against strong solver and heuristic baselines. With an evolving space, give policies access to the same transformation and component libraries, then measure which designs they expose and how much evaluation effort they spend. A retention ablation should revisit the conditional reversals seen in the case studies.

Knowledge transfer is also conditional. A mapping rule learned on one architecture can identify a promising mechanism elsewhere, but its performance threshold depends on the new resource ratios. The negative transfer case argues for storing those conditions and validating them on reserved workloads. Human intervention and model-inference costs belong in the accounting. Reusable formulation knowledge is valuable when its construction cost is amortized over enough related design problems.

\section{Conclusion}
AI infrastructure design joins algorithmic choices, programs, interfaces, and hardware into one coupled optimization problem. Its difficulty lies both in representing valuable implementations and in searching the resulting combinations. Agentic Progressive Co-Design addresses these tasks together: introduce the mechanisms needed to express a useful design, encode their interactions in a MIP, exploit structure to make search tractable, and use paired implementation feedback to revise the next model.

The Ventus studies show that hardware changes can reverse software preferences, that changing both sides improves over optimizing each alone, and that a closed solver gap cannot repair an inaccurate performance model. These findings motivate a design process in which abstraction boundaries and optimization formulations remain revisable. The resulting research direction is a systematic route from workload semantics to a matched algorithm, executable program, software--hardware interface, and hardware organization.
